\textbf{Pour cette première partie, aucune justification n'est demandée et une seule réponse est possible par question. Pour chaque question, reportez son numéro sur votre copie et indiquez votre réponse.}

Une réponse fausse ou l'absence de réponse n'enlève aucun point.

\subsection*{Question 1}
	On note $S$ l'ensemble des solutions dans $\R$ de l'équation : \quad $(2 x-6)(x+4)=0$. \quad On a :

\medskip
\begin{tblr}{colspec = {*4{X}}}
\textbf{A.}\quad $S=\{4 ; 6\}$
&
\textbf{B.}\quad $S=\{-4 ; 6\}$
&
\textbf{C.}\quad $S=\{3 ; 4\}$
&
\textbf{D.}\quad $S=\{-4 ; 3\}$
\\
\end{tblr}

\subsection*{Question 2}
	On lance un dé cubique équilibré numéroté de 1 à 6. Quelle est la probabilité d'obtenir un nombre strictement supérieur à 4 ?

\medskip
\begin{tblr}{colspec = {*4{X}}}
\textbf{A.}\quad $\dfrac{1}{2}$
&
\textbf{B.}\quad $\dfrac{1}{3}$
&
\textbf{C.}\quad $\dfrac{1}{4}$
&
\textbf{D.}\quad $\dfrac{1}{5}$
\\
\end{tblr}

\subsection*{Question 3}
	La seule égalité vraie est :

\medskip
\begin{tblr}{colspec = {*4{X}}}
\textbf{A.}\quad $3^{-5} \times 6^{-5}=18^{-10}$
&
\textbf{B.}\quad $\left(3^{-5}\right)^{2}=3^{-3}$
&
\textbf{C.}\quad $\dfrac{3}{3^{5}}=3^{-4}$
&
\textbf{D.}\quad $\dfrac{3^{-2}}{3^{8}}=3^{-6}$
\\
\end{tblr}

\needspace{10\baselineskip}
\subsection*{Question 4}
	On a représenté ci-dessous quatre droites dans un repère orthonormé.

	\begin{center}
	\begin{tikzpicture}[x =0.9cm,y=0.9cm,>=stealth]
		\draw[gray!75,line width=0.4pt,xstep=1, ystep=1] (-8.4,-4.5) grid (8.2,7.9) ;
		\draw[->, line width=1.2pt] (-8.4,0)--(8.2,0) ;
		\draw[->, line width=1.2pt] (0,-4.5)--(0,7.9) ;
		\foreach \x in {-8,...,-1,1,2,...,8}
			\draw[thick] (\x,2pt)--(\x,-2pt) node[below]{$\np{\x}$} ;
		\foreach \y in {-4,...,-1,1,2,...,7}
			\draw[thick] (2pt,\y)--(-2pt,\y) node[left]{$\np{\y}$} ;
		\node[below left] at (-2pt,-2pt) {0};
		\draw[line width=1.4pt,red] (3,7.9)--(3,-4.5) node[pos=0.1,left] {$\mathrm{D}_1$} ;
		\draw[line width=1.4pt,teal] (-7.9,7.9)--(4.5,-4.5) node[pos=0.15,above right] {$\mathrm{D}_2$} ;
		\draw[line width=1.4pt,blue5] (-8.4,3)--(8.2,3) node[pos=0.05,above] {$\mathrm{D}_3$} ;
		\draw[line width=1.4pt,purple5] (-5.5,-4.5)--(6.9,7.9) node[pos=0.11,above left] {$\mathrm{D}_4$} ;
	\end{tikzpicture}
	\end{center}

	Parmi les quatre affirmations suivantes, laquelle est vraie ?

\medskip
\begin{tblr}{colspec = {*2{X}}}
\textbf{A.}\quad La droite $\mathrm{D}_{1}$ a pour équation $y=3$
&
\textbf{B.}\quad La droite $\mathrm{D}_{2}$ a pour coefficient directeur 1
\\
\textbf{C.}\quad La droite $\mathrm{D}_{3}$ a pour coefficient directeur 3
&
\textbf{D.}\quad La droite $\mathrm{D}_{4}$ a pour équation $y=x+1$
\\
\end{tblr}

\subsection*{Question 5}
	Soit $\mathcal{C}_{f}$ la courbe représentative de la fonction $f$ définie sur $\R$ par \quad $f(x)=2 x^{2}-3 x+1$.

	Parmi les quatre points suivants, lequel appartient à $\mathcal{C}_{f}$ ?

\medskip
\begin{tblr}{colspec = {*4{X}}}
\textbf{A.}\quad E$(-1 ; 2)$
&
\textbf{B.}\quad H$(-3 ; 28)$
&
\textbf{C.}\quad G$(2 ; 1)$
&
\textbf{D.}\quad K$(3 ; 4)$
\\
\end{tblr}

\subsection*{Question 6}
	On considère le nombre \quad $E=\left(\dfrac{1}{2}+\dfrac{1}{3}\right) \times \dfrac{2}{5}$. \quad Une écriture simplifiée de $E$ est :

\medskip
\begin{tblr}{colspec = {*4{X}}}
\textbf{A.}\quad $\dfrac{4}{25}$
&
\textbf{B.}\quad $\dfrac{1}{3}$
&
\textbf{C.}\quad $\dfrac{2}{30}$
&
\textbf{D.}\quad $\dfrac{4}{5}$
\\
\end{tblr}

\needspace{5\baselineskip}
\subsection*{Question 7}
	Le prix d'un pantalon a augmenté de $10\,\%$ puis a baissé de $20\,\%$.

	Quelle évolution globale a subi le prix du pantalon ?

\medskip
\begin{tblr}{colspec = {*2{X}}}
\textbf{A.}\quad une baisse de 15\,\%
&
\textbf{B.}\quad une augmentation de 8,8\,\%
\\
\textbf{C.}\quad une baisse de 12\,\%
&
\textbf{D.}\quad une baisse de 10\,\%
\\
\end{tblr}

\subsection*{Question 8}
	L'accélération gravitationnelle $g$ (en $\mathrm{m} / \mathrm{s}^{2}$) d'un astre est donnée par la formule
	\[g=\dfrac{G m}{r^{2}}\]
	où $G$ est la constante gravitationnelle universelle (en $\mathrm{N}\cdot \mathrm{m}^{2} / \mathrm{kg}^{2}$),\\
	$m$ est la masse de l'astre (en kg) et $r$ le rayon de l'astre (en m).

	L'expression permettant, à partir de cette formule, d'exprimer le rayon $r$ est :

\medskip
\begin{tblr}{colspec = {*4{X}}}
\textbf{A.}\quad $ r=\dfrac{G m}{g^{2}} $
&
\textbf{B.}\quad $ r=\sqrt{G m g} $
&
\textbf{C.}\quad $ r=\sqrt{\dfrac{G m}{g}} $
&
\textbf{D.}\quad $ r=\sqrt{\dfrac{g}{G m}} $
\\
\end{tblr}

\needspace{10\baselineskip}
